\chapter{Front End Implementation}
\label{chap:front-end}

The front end of a compiler converts source text into a structured form that can
be analyzed. In \epl{}, the front end consists of the lexer and parser. The
lexer produces tokens with spans. The parser consumes those tokens and constructs
an AST. This chapter explains the implementation choices and how the source
language is represented before semantic analysis.

\section{Lexer Design}

The lexer is implemented in \texttt{src/lex.rs}. It is a Rust iterator over
\texttt{Result<(Span, Token), Error>}. Each successful item contains a span and a
token. Each failed item contains a lexical error with the span where the problem
was found.

The lexer stores two pieces of state: the current byte offset and a reference to
the source string. On each iteration it skips whitespace and comments, then tries
to recognize the next token. Comments start with \texttt{\#} and continue until
the end of the line. This simple rule avoids interaction with nested block
comments or preprocessor-like behavior.
Figure~\ref{fig:lexer-flow} shows the control flow of this token-reading loop.

\begin{figure}[H]
  \centering
  \fbox{\begin{minipage}{0.86\textwidth}
  \centering
  \begin{tabular}{c}
  read next character \\
  $\downarrow$ \\
  whitespace or comment? skip and repeat \\
  $\downarrow$ \\
  identifier start? read identifier or keyword \\
  $\downarrow$ \\
  digit? read number and optional suffix \\
  $\downarrow$ \\
  quote? read string literal \\
  $\downarrow$ \\
  punctuation prefix? emit punctuation \\
  $\downarrow$ \\
  otherwise lexical error
  \end{tabular}
  \end{minipage}}
  \caption{Control flow of the \epl{} lexer}
  \label{fig:lexer-flow}
\end{figure}

\section{Tokens}

The token enum has four main variants: keyword, identifier, literal, and
punctuation. Keywords are stored in a fixed keyword enum. Identifiers store their
source string. Literals currently include number literals and string literals.
Punctuation includes operators and delimiters such as \texttt{+},
\texttt{==}, \texttt{->}, \texttt{...}, the dot-left-brace token, parentheses, braces, and
brackets.

The lexer checks punctuation by trying entries in a punctuation table ordered
from longer strings to shorter strings. This matters because \texttt{...} must
be recognized before \texttt{..}, \texttt{==} before \texttt{=}, and
\texttt{->} before \texttt{-}. The implementation keeps this priority visible in
one table.

Number literals are parsed as \texttt{i128} first. This lets the lexer accept a
large neutral representation before semantic analysis chooses the target integer
type. If the number text does not fit into \texttt{i128}, the lexer reports
\texttt{NumberLiteralTooLarge}. A number may be followed by an identifier suffix
without whitespace, such as \texttt{42u64}. The suffix is stored with its own
span so semantic analysis can report unknown suffixes precisely.

String literals support a small set of escape sequences for newline, carriage
return, tab, and a literal backslash. Unknown escape sequences and unclosed
strings are lexical errors. The language does not currently include character
literals or raw strings.

\section{Spans}

A span is a half-open byte range \texttt{start..end} into the source string.
Spans are attached to tokens, AST identifiers, type nodes, expression nodes, and
many semantic errors. The \texttt{Span::join} method creates a larger span
covering two smaller spans. This is useful when an expression spans multiple
tokens, for example when a binary expression spans from the start of its left
operand to the end of its right operand.

The diagnostic printer uses spans to compute the line number and underline. This
is why even small AST nodes, such as identifiers and punctuation-driven
constructs, preserve source locations. A compiler that drops source locations in
the parser can still compile code, but it cannot explain errors well.

\section{Parser Design}

The parser is implemented in \texttt{src/ast.rs}. It owns the lexer and a small
lookahead queue. The parser exposes a single public entry point:
\texttt{Parser::new(src).parse()}. The result is an \texttt{Ast} containing a
vector of top-level items.

The parser is recursive descent. Each grammar category is represented by a
method. For example, \texttt{next\_function} parses a function item,
\texttt{next\_type} parses a type, \texttt{next\_block\_expr} parses a block,
and expression precedence is implemented by methods such as
\texttt{next\_or\_expr}, \texttt{next\_and\_expr},
\texttt{next\_comp\_expr}, \texttt{next\_additive\_expr}, and
\texttt{next\_multiplicative\_expr}.

The parser performs syntax validation but not full semantic validation. It can
reject duplicate annotations on an item, reject \texttt{let name;} without type
or initializer, require variadic \texttt{...} to be the final function
parameter, and report unexpected tokens. It does not check whether a type name
exists, whether a variable is in scope, whether a function call has the correct
argument types, or whether an arithmetic operator is applied to integers. Those
checks belong to semantic lowering.

\section{Top-Level Items}

An \epl{} source file is parsed as zero or more top-level items. Each item may
have annotations. The parser stores annotations in a \texttt{BTreeSet} keyed by
annotation name, which makes duplicate detection simple.

Function parsing handles both definitions and declarations. The parser reads the
\texttt{fn} keyword, the function name, parameter list, optional return type,
and then either a block body or a semicolon. A semicolon means that the function
is a declaration. Variadic functions are represented by a boolean flag on the
AST function.

Struct parsing reads a name and a comma-delimited list of fields. Field types
are parsed as AST type nodes, not resolved to semantic types yet. This is
important because struct processing later uses the global type namespace and
computes layout.

\section{Expression Parsing}

The expression grammar is layered by precedence. Assignment has low precedence,
logical operators are above assignment, comparison is above logical operators,
range is above comparison, addition is above range, multiplication is above
addition, casts are above multiplication, unary operators are above casts, and
field/index/dereference access is above unary operations.

The parser treats several constructs as expression forms:
\texttt{return}, \texttt{break}, \texttt{continue}, \texttt{comptime}, blocks,
\texttt{if}, \texttt{loop}, \texttt{while}, \texttt{for}, array initializers,
struct initializers, calls, casts, unary operations, binary operations,
assignments, identifiers, and literals.

Blocks require a small parsing rule because \epl{} distinguishes statements
from the final expression. A block may contain semicolons, \texttt{let}
statements, expression statements, and an optional final expression. If the next
expression is followed by a closing brace, it becomes the final expression. If it
is followed by a semicolon, it becomes an expression statement.

\section{AST Shape}

The AST is deliberately close to source syntax. The \texttt{Expr} enum includes
variants for \texttt{While}, \texttt{For}, \texttt{Range}, \texttt{LogicalAnd}
and \texttt{LogicalOr} as binary operations, \texttt{AsCast}, \texttt{Comptime},
and other source-level forms. Later stages lower these into a smaller semantic
core. For example, \texttt{while} does not survive as a distinct lower construct;
it becomes a \texttt{loop} containing an \texttt{if}.

Keeping the AST close to source syntax has two advantages. First, it makes AST
debug output useful when diagnosing parser behavior. Second, it avoids mixing
syntax parsing with language design decisions that are easier to make after
types are known.

\section{Parser Error Handling}

Parser errors contain an optional span and an error kind. Unexpected token errors
record both the expected description and the token that was actually found.
Lexical errors are wrapped so the parser can expose one error type to the rest
of the compiler.

The parser stops at the first error. It does not implement error recovery or
multi-error reporting. For a teaching compiler this is acceptable, but it is a
clear area for future improvement. Production compilers often try to recover
after syntax errors so they can report several problems in one run.
